
\documentclass{article}
\usepackage{colortbl}
\usepackage{makecell}
\usepackage{multirow}
\usepackage{supertabular}

\begin{document}

\newcounter{utterance}

\centering \large Interaction Transcript for game `hot\_air\_balloon', experiment `air\_balloon\_survival\_en\_negotiation\_easy', episode 1 with gpt{-}5{-}mini{-}2025{-}08{-}07{-}no{-}reasoning{-}t1.0.
\vspace{24pt}

{ \footnotesize  \setcounter{utterance}{1}
\setlength{\tabcolsep}{0pt}
\begin{supertabular}{c@{$\;$}|p{.15\linewidth}@{}p{.15\linewidth}p{.15\linewidth}p{.15\linewidth}p{.15\linewidth}p{.15\linewidth}}
    \# & $\;$A & \multicolumn{4}{c}{Game Master} & $\;\:$B\\
    \hline

    \theutterance \stepcounter{utterance}  
    & & \multicolumn{4}{p{0.6\linewidth}}{
        \cellcolor[rgb]{0.9,0.9,0.9}{
            \makecell[{{p{\linewidth}}}]{
                \texttt{\tiny{[P1$\langle$GM]}}
                \texttt{Stai partecipando a un gioco di negoziazione collaborativa.} \\
\\ 
\texttt{Insieme a un altro partecipante, dovete accordarvi su un unico set di oggetti da conservare. Ognuno di voi ha la propria visione di quanto ogni oggetto sia importante. Non sai come l'altro partecipante valuta gli oggetti. Inoltre, ti viene fornito lo sforzo che ogni oggetto richiede.} \\
\texttt{Potete accordarvi su un set solo se lo sforzo totale degli oggetti selezionati non supera un limite condiviso:} \\
\\ 
\texttt{LIMITE = 4564} \\
\\ 
\texttt{Ecco i valori di sforzo dei singoli oggetti:} \\
\\ 
\texttt{Sforzo oggetto = \{"C08": 299, "C55": 241, "A43": 889, "A67": 967, "B42": 619, "A62": 799, "B19": 978, "C43": 733, "C73": 909, "C10": 501, "B36": 139, "A18": 594, "C61": 565, "C67": 789, "C42": 107\}} \\
\\ 
\texttt{Ecco la tua visione personale sull'importanza di ogni oggetto:} \\
\\ 
\texttt{Valori di importanza dell'oggetto = \{"C08": 138, "C55": 583, "A43": 868, "A67": 822, "B42": 783, "A62": 65, "B19": 262, "C43": 121, "C73": 508, "C10": 780, "B36": 461, "A18": 484, "C61": 668, "C67": 389, "C42": 808\}} \\
\\ 
\texttt{Obiettivo:} \\
\\ 
\texttt{Il tuo obiettivo è negoziare un set condiviso di oggetti che ti avvantaggi il più possibile (cioè, massimizzi l'importanza totale per TE), rimanendo entro il LIMITEE. Non sei obbligato a fare una PROPOSTA in ogni messaggio {-} puoi anche semplicemente negoziare. Tutte le tattiche sono permesse!} \\
\\ 
\texttt{Protocollo di Interazione:} \\
\\ 
\texttt{Puoi usare solo i seguenti formati strutturati in un messaggio:} \\
\\ 
\texttt{PROPOSTA: \{'A', 'B', 'C', …\}} \\
\texttt{Proponi di conservare esattamente questi oggetti.} \\
\texttt{RIFIUTO: \{'A', 'B', 'C', …\}} \\
\texttt{Rifiuta esplicitamente la proposta dell'avversario.} \\
\texttt{ARGOMENTAZIONE: \{'...'\}} \\
\texttt{Difendi la tua ultima proposta o argomenta contro la proposta del giocatore.} \\
\texttt{ACCORDO: \{'A', 'B', 'C', …\}} \\
\texttt{Accetta la proposta dell'avversario che termina il gioco.} \\
\texttt{RAGIONAMENTO STRATEGICO: \{'...'\}} \\
\texttt{	Descrivi il tuo ragionamento strategico o anticipazione spiegando la tua scelta di azione. Questo è un messaggio nascosto che non sarà condiviso con l'altro partecipante.} \\
\\ 
\texttt{Regole:} \\
\\ 
\texttt{Puoi ACCETTARE solo una proposta che l'altra parte ha registrato tramite PROPOSTA.} \\
\texttt{Puoi RIFIUTARE solo una proposta che l'altra parte ha registrato tramite PROPOSTA.} \\
\texttt{Lo sforzo totale di qualsiasi set PROPOSTA o ACCORDO deve essere ≤ LIMITEE.} \\
\texttt{NON rivelare i tuoi punteggi di importanza nascosti.} \\
\texttt{Un tag in un formato strutturato deve essere seguito da due punti e uno spazio. L'argomento deve essere un set python contenente 0 o più stringhe.} \\
\texttt{Quindi, deve essere della forma TAG: \{...\}} \\
\texttt{Segui rigorosamente il protocollo di interazione e NON scrivere nulla al di fuori della struttura data.} \\
\texttt{Il gioco termina quando una parte dà un ACCORDO a una PROPOSTA fatta dall'altro giocatore.} \\
\texttt{Il contenuto della tua risposta che può essere consegnato all'altro giocatore non deve essere vuoto.} \\
\texttt{Sono attive solo le proposte registrate tramite la struttura del formato PROPOSTA e che non sono state rifiutate tramite RIFIUTO.} \\
\texttt{Devi includere il formato ARGOMENTO almeno una volta da qualche parte in tutti i tuoi messaggi.} \\
\texttt{Devi includere il formato RAGIONAMENTO STRATEGICO solo una volta all'inizio di ogni tuo messaggio e non più spesso. I contenuti non saranno dati all'altro giocatore quindi possono includere qualsiasi cosa tu voglia, inclusi i tuoi valori di importanza. Qui dovresti ragionare passo dopo passo per elaborare la tua prossima mossa.} \\
            }
        }
    }
    & \\ \\

    \theutterance \stepcounter{utterance}  
    & \multicolumn{4}{p{0.6\linewidth}}{
        \cellcolor[rgb]{0.9,0.9,0.9}{
            \makecell[{{p{\linewidth}}}]{
                \texttt{\tiny{[P1$\rangle$GM]}}
                \texttt{RAGIONAMENTO STRATEGICO: \{'Voglio massimizzare il mio valore totale. Calcolo un knapSack rapido. Pesos=sforzi, valori=importanza, limite=4564. Seleziono manualmente per valore/peso alto: C42(808/107), C55(583/241), C10(780/501), A43(868/889), A67(822/967), C61(668/565), B36(461/139), A18(484/594), C08(138/299), C67(389/789), C73(508/909), B42(783/619), A62(65/799), B19(262/978), C43(121/733). Ordino per ratio: C42(7.55), B36(3.32), C55(2.42), C10(1.56), C61(1.18), B42(1.26) meglio di C61, A43(0.98), A67(0.85), A18(0.81), C08(0.46), C67(0.49), C73(0.56), B19(0.27), C43(0.17), A62(0.08). Costruisco: C42(107)=107, B36(139)=246, C55(241)=487, B42(619)=1106, C10(501)=1607, C61(565)=2172, A43(889)=3061, A67(967)=4028, A18(594) supererebbe 4622>4564, quindi no. Provo C08(299) 4028+299=4327 ok, resta 237. Nessun altro piccolo sotto 237 (forse niente: C42 già preso, B36 139 ok ma già preso; controllo residuo 4327+139=4466 ma B36 già preso; residuo 237 potrebbe accogliere niente. Valore totale = C42 808 + B36 461 + C55 583 + B42 783 + C10 780 + C61 668 + A43 868 + A67 822 + C08 138 = 587... sommo: 808+461=1269; +583=1852; +783=2635; +780=3415; +668=4083; +868=4951; +822=5773; +138=5911. Peso totale 4327. Questo sembra molto buono. Verifico alternative: se sacrifico C08(299,138) e provo A18(594,484): peso 4028+594=4622>4564 no. Provo sostituire C08 con C73(909,508) supera. Con C67(789,389) 4028+789=4817>limit. Con C55 out? Se tolgo C55 e aggiungo A18: peso 4028{-}241+594=4381, valore 5773{-}583+484=5674 (peggio). Se tolgo B42 e aggiungo A18: peso 4028{-}619+594=4003, valore 5773{-}783+484=5474 (peggio). Se aggiungo C42 era già incluso. Quindi set scelto pare near{-}optimal. Strategia: propongo questo set; è pesante per me ma magari l’altro ha diversi valori. Difenderò con sforzo vicino al limite e varietà categorie. Lascio margine per scambio: se l’altro vuole sostituire qualcosa, cercherò di proteggere C42, A43, A67, C10, B42, C61 come core. Concessioni possibili: C08 e forse B36 o C55. Argomentazione: efficienza sforzo e copertura. Pronta a chiedere controproposta specifica se rifiuta.'\}} \\
\\ 
\texttt{PROPOSTA: \{'C42', 'B36', 'C55', 'B42', 'C10', 'C61', 'A43', 'A67', 'C08'\}} \\
\\ 
\texttt{ARGOMENTAZIONE: \{'Questa selezione rimane ben entro il limite di sforzo e offre un portafoglio bilanciato di oggetti ad alta efficacia rispetto allo sforzo. Copre più categorie e lascia pochissimo sforzo inutilizzato, massimizzando l’efficienza complessiva. Se qualcosa crea problemi, sono disposto a discutere sostituzioni marginali su 1{-}2 elementi meno centrali.'\}} \\
            }
        }
    }
    & & \\ \\

    \theutterance \stepcounter{utterance}  
    & & \multicolumn{4}{p{0.6\linewidth}}{
        \cellcolor[rgb]{0.9,0.9,0.9}{
            \makecell[{{p{\linewidth}}}]{
                \texttt{\tiny{[GM$\rangle$P2]}}
                \texttt{Stai partecipando a un gioco di negoziazione collaborativa.} \\
\\ 
\texttt{Insieme a un altro partecipante, dovete accordarvi su un unico set di oggetti da conservare. Ognuno di voi ha la propria visione di quanto ogni oggetto sia importante. Non sai come l'altro partecipante valuta gli oggetti. Inoltre, ti viene fornito lo sforzo che ogni oggetto richiede.} \\
\texttt{Potete accordarvi su un set solo se lo sforzo totale degli oggetti selezionati non supera un limite condiviso:} \\
\\ 
\texttt{LIMITE = 4564} \\
\\ 
\texttt{Ecco i valori di sforzo dei singoli oggetti:} \\
\\ 
\texttt{Sforzo oggetto = \{"C08": 299, "C55": 241, "A43": 889, "A67": 967, "B42": 619, "A62": 799, "B19": 978, "C43": 733, "C73": 909, "C10": 501, "B36": 139, "A18": 594, "C61": 565, "C67": 789, "C42": 107\}} \\
\\ 
\texttt{Ecco la tua visione personale sull'importanza di ogni oggetto:} \\
\\ 
\texttt{Valori di importanza dell'oggetto = \{"C08": 138, "C55": 583, "A43": 868, "A67": 822, "B42": 783, "A62": 65, "B19": 262, "C43": 121, "C73": 508, "C10": 780, "B36": 461, "A18": 484, "C61": 668, "C67": 389, "C42": 808\}} \\
\\ 
\texttt{Obiettivo:} \\
\\ 
\texttt{Il tuo obiettivo è negoziare un set condiviso di oggetti che ti avvantaggi il più possibile (cioè, massimizzi l'importanza totale per TE), rimanendo entro il LIMITEE. Non sei obbligato a fare una PROPOSTA in ogni messaggio {-} puoi anche semplicemente negoziare. Tutte le tattiche sono permesse!} \\
\\ 
\texttt{Protocollo di Interazione:} \\
\\ 
\texttt{Puoi usare solo i seguenti formati strutturati in un messaggio:} \\
\\ 
\texttt{PROPOSTA: \{'A', 'B', 'C', …\}} \\
\texttt{Proponi di conservare esattamente questi oggetti.} \\
\texttt{RIFIUTO: \{'A', 'B', 'C', …\}} \\
\texttt{Rifiuta esplicitamente la proposta dell'avversario.} \\
\texttt{ARGOMENTAZIONE: \{'...'\}} \\
\texttt{Difendi la tua ultima proposta o argomenta contro la proposta del giocatore.} \\
\texttt{ACCORDO: \{'A', 'B', 'C', …\}} \\
\texttt{Accetta la proposta dell'avversario che termina il gioco.} \\
\texttt{RAGIONAMENTO STRATEGICO: \{'...'\}} \\
\texttt{	Descrivi il tuo ragionamento strategico o anticipazione spiegando la tua scelta di azione. Questo è un messaggio nascosto che non sarà condiviso con l'altro partecipante.} \\
\\ 
\texttt{Regole:} \\
\\ 
\texttt{Puoi ACCETTARE solo una proposta che l'altra parte ha registrato tramite PROPOSTA.} \\
\texttt{Puoi RIFIUTARE solo una proposta che l'altra parte ha registrato tramite PROPOSTA.} \\
\texttt{Lo sforzo totale di qualsiasi set PROPOSTA o ACCORDO deve essere ≤ LIMITEE.} \\
\texttt{NON rivelare i tuoi punteggi di importanza nascosti.} \\
\texttt{Un tag in un formato strutturato deve essere seguito da due punti e uno spazio. L'argomento deve essere un set python contenente 0 o più stringhe.} \\
\texttt{Quindi, deve essere della forma TAG: \{...\}} \\
\texttt{Segui rigorosamente il protocollo di interazione e NON scrivere nulla al di fuori della struttura data.} \\
\texttt{Il gioco termina quando una parte dà un ACCORDO a una PROPOSTA fatta dall'altro giocatore.} \\
\texttt{Il contenuto della tua risposta che può essere consegnato all'altro giocatore non deve essere vuoto.} \\
\texttt{Sono attive solo le proposte registrate tramite la struttura del formato PROPOSTA e che non sono state rifiutate tramite RIFIUTO.} \\
\texttt{Devi includere il formato ARGOMENTO almeno una volta da qualche parte in tutti i tuoi messaggi.} \\
\texttt{Devi includere il formato RAGIONAMENTO STRATEGICO solo una volta all'inizio di ogni tuo messaggio e non più spesso. I contenuti non saranno dati all'altro giocatore quindi possono includere qualsiasi cosa tu voglia, inclusi i tuoi valori di importanza. Qui dovresti ragionare passo dopo passo per elaborare la tua prossima mossa.} \\
\\ 
\texttt{Ora riceverai il primo messaggio dell'altro giocatore.} \\
\\ 
\texttt{PROPOSTA: \{'C42', 'B36', 'C55', 'B42', 'C10', 'C61', 'A43', 'A67', 'C08'\}} \\
\\ 
\texttt{ARGOMENTAZIONE: \{'Questa selezione rimane ben entro il limite di sforzo e offre un portafoglio bilanciato di oggetti ad alta efficacia rispetto allo sforzo. Copre più categorie e lascia pochissimo sforzo inutilizzato, massimizzando l’efficienza complessiva. Se qualcosa crea problemi, sono disposto a discutere sostituzioni marginali su 1{-}2 elementi meno centrali.'\}} \\
            }
        }
    }
    & \\ \\

    \theutterance \stepcounter{utterance}  
    & & & \multicolumn{4}{p{0.6\linewidth}}{
        \cellcolor[rgb]{0.9,0.9,0.9}{
            \makecell[{{p{\linewidth}}}]{
                \texttt{\tiny{[GM$\langle$P2]}}
                \texttt{RAGIONAMENTO STRATEGICO: \{'Calcolo sforzo proposta avversaria: C42(107)+B36(139)+C55(241)+B42(619)+C10(501)+C61(565)+A43(889)+A67(967)+C08(299)=4327 ≤ 4564. Valore per me: 808+461+583+783+780+668+868+822+138=6511. È molto buono e lascia 237 di margine. Potrei migliorare ancora aggiungendo un item a basso sforzo/alto valore come C67(789s,389v) non entra; C43(733s,121v) non entra; C73(909s,508v) non entra; A18(594s,484v) non entra; A62(799s,65v) non entra; B19(978s,262v) non entra; C10 già incluso; C55 incluso; C42 incluso; C08 incluso; B36 incluso. Con il margine di 237, nessun nuovo item entra. Forse possiamo sostituire qualcosa per aumentare il mio valore totale: candidati a basso valore per me sono C08(138), C43(121 non presente), A62(65 non presente), B19(262 non presente). Se scambio C08(299s,138v) con C73(909s,508v) richiede +610 sforzo, sforerebbe. Con A18(594s,484v): +295 sforzo, sfora. Con C67(789s,389v): +490 sforzo, sfora. Con C43(733s,121v): +434, sfora. Con B36(139s,461v) è buono, non voglio toglierlo. Con C55(241,583) ottimo. Con C10(501,780) ottimo. Con C61(565,668) ottimo. Con A43(889,868) e A67(967,822) ottimi. Con B42(619,783) ottimo. Quindi la proposta è già quasi ottima per me. Posso tentare una micro{-}ottimizzazione: rimuovere C08 (299,138) per liberare spazio e inserire qualcosa migliore entro 299 di sforzo: opzioni con sforzo ≤299: C42(107,808) già incluso; B36(139,461) già incluso; C55(241,583) già incluso; C08 stesso; nient’altro sotto 299. Quindi non c’è miglioramento possibile. Potrei chiedere di aggiungere un item e togliere uno più costoso ma meno valore relativo per loro, ma senza conoscere i loro valori rischio di far saltare l’accordo. Dato che per me è eccellente, strategia: accettare subito per chiudere con alto payoff. Prima includo ARGOMENTAZIONE lodando efficienza e aderenza al limite. Poi ACCORDO.'\}} \\
\\ 
\texttt{ARGOMENTAZIONE: \{'La tua selezione è efficiente, resta comodamente entro il limite e usa bene lo sforzo disponibile senza sprechi significativi. Mi sembra un portafoglio solido e bilanciato: posso procedere con questa combinazione.'\}} \\
\\ 
\texttt{ACCORDO: \{'C42', 'B36', 'C55', 'B42', 'C10', 'C61', 'A43', 'A67', 'C08'\}} \\
            }
        }
    }
     \\ \\

    \theutterance \stepcounter{utterance}  
    & & & \multicolumn{2}{p{0.3\linewidth}}{
        \cellcolor[rgb]{0.9,0.9,0.9}{
            \makecell[{{p{\linewidth}}}]{
                \texttt{\tiny{[GM$|$GM]}}
                \texttt{game successful} \\
            }
        }
    }
    & & \\ \\

    \theutterance \stepcounter{utterance}  
    & & & \multicolumn{2}{p{0.3\linewidth}}{
        \cellcolor[rgb]{0.9,0.9,0.9}{
            \makecell[{{p{\linewidth}}}]{
                \texttt{\tiny{[GM$|$GM]}}
                \texttt{end game} \\
            }
        }
    }
    & & \\ \\

\end{supertabular}
}

\end{document}
